\chapter{Statistical and Analytical Strategy\label{cha:chapter4}}

This chapter details the analytical methodology, statistical models, and hypothesis testing strategies employed to evaluate experimental results.

\section{Primary Contrasts\label{sec:primary-contrasts}}

To analyze prompt factors, we compare each experimental condition against a standard baseline: neutral tone, no role, and plain text. Correctness differences are expressed in percentage points of the success rates, whereas behavioral effects are expressed in the number of tokens used and response duration. The results are supplemented by p-values; however, they are considered exploratory unless adjusted for multiple hypothesis testing.

The use of percentage-point differences is advantageous since they provide an easily understandable basis for prompt-engineering decisions. For instance, a shift in the pass rate of four points corresponds to four additional or fewer successes out of 100 trials, without having to account for cost and latency. It simplifies the engineering heuristic compared to odds ratios. For this reason, we focus on the effect sizes and their directions, and we regard p-values merely as measures of statistical uncertainty rather than binary thresholds for significance.

\section{Stratification by Task Kind and Model Group}

We have separated our results according to task domain (code vs. math) and deployment method (local vs. API). This stratification is critical because our tests show that the effect of the prompt varies depending on both the task type and the model used. Using data from just one cumulative measure can cause us to miss important variations. For example, JSON and XML formats improve performance on API-enabled Python tasks; conversely, for mathematical tasks run locally, the use of JSON formats decreases accuracy.

In addition, this stratification provides protection from overgeneralization. A pooled estimate may be statistically precise but risks leading to misguided conclusions if it averages together different mechanisms. Code generation relies on executable artifacts and parser compatibility, while mathematical reasoning depends on answer extraction and numerical evaluation. Local and API models differ in their default verbosity, instruction adherence, latency, and serving infrastructure. Rather than treating these variations as statistical noise to be averaged out, our analysis considers them integral components of the phenomenon under investigation.

\section{Mixed-Effects Models as the Main Inferential Basis}

As each base task is assessed multiple times according to all 18 prompt types and under each model in the stratum, the pass/fail observations are not independent. Mixed-effects logistic regression is the appropriate inferential framework for such clustered benchmark data \cite{gelman2007data}. Therefore, the primary analysis in this thesis is a mixed-effects logistic regression for pass/fail status carried out separately for local code, local math, API Python, and API math. Each model includes fixed effects for the five contrasts of prompt factors (aggressive and polite vs. neutral tone, role framing vs. no role, and JSON and XML vs. plaintext), fixed effects for models in the stratum, and a random intercept per base task with a single shared variance component, which allows modeling the repetition of the same 50 base tasks explicitly instead of treating them as independent rows. Model estimation is done using variational Bayes as implemented in statsmodels' \texttt{BinomialBayesMixedGLM}. We provide the results in the form of odds ratios with 95\% intervals and average marginal effects (AME) in percentage points, averaged over the observed covariates and the estimated task random-intercept distribution, with intervals from posterior draws.

Similarly, the repeated-measures structure is applied to the behavioral outcomes. Log completion tokens and log duration are modeled with linear mixed-effects models using the same fixed-effect structure and a base-task random intercept. Result-only compliance, being binary, is modeled with the same logistic mixed-effects specification as pass/fail. However, for API code, compliance remains constant under the system prompts, precluding model estimation; therefore, these results are reported descriptively.

The percentage-point contrasts and linear probability models (LPMs) with model-level and base-task fixed effects are retained as a descriptive and interpretable layer: a 4 percentage-point increase means 4 more (or fewer) successful trials out of 100, which is easier to handle from an engineering perspective than odds ratios. Furthermore, linear probability models with standard errors clustered by the base task are kept as a robustness check against mixed-effects estimates. We note that predicted probabilities from linear probability models are not bounded by the unit interval (i.e., they can fall outside the 0 to 1 range); therefore, clustered LPMs are not seen as formal estimates but rather as sensitivity measures, while logistic regression remains the main method of inference for pass/fail claims.
Robustness expectation is directional rather than exact. P-values for the mixed model, the clustered LPM, and raw contrasts will not necessarily be identical, but the sign and approximate magnitude of the headline effects should agree. Where they do not, we investigate and report the discrepancy rather than presenting an isolated significant result as a headline finding.

\section{Multiple Comparisons Correction}

Given the many different comparisons made in this experiment, both in terms of factors and models, p-values must be interpreted carefully. We report uncorrected p-values on the main contrast tables for readability; subsequently, we apply the Benjamini-Hochberg false discovery rate correction to evaluate model-specific local baseline contrasts. Findings are considered provisional if they fall outside the corrected list of model-specific baselines, such as the impacts of API math role framing.

The goal of the correction approach is intentionally targeted. This thesis makes no claims or assumptions about all possible pairs of comparisons being confirmatory. Instead, it separates headline aggregate contrasts, corrected local model-based baseline contrasts, and exploratory details to ensure that the story does not turn isolated low p-values into broad prompt rules. Findings that survive correction or have been determined to be directionally consistent across robustness testing will receive greater weight than isolated marginal contrasts.

This distinction shapes the reporting language throughout the thesis. A finding can have interesting practical significance even when it is not treated as confirmatory, especially if it highlights a possible implementation issue such as higher token costs or lower compliance. On the other hand, a finding that is statistically significant but has a small effect size may carry limited importance in engineering. The interpretation of findings is therefore dependent on effect size, uncertainty, robustness, and operational relevance.

\section{Mediation and Behavioral Controls}

The mediation-style analysis uses behavioral decomposition rather than proving the existence of hidden reasoning processes. Log completion tokens, result-only compliance, and log duration are the adjustment variables. The goal is to describe how much of each observed prompt contrast remains after adjustment for observable behavior related to cost, latency, and parser compatibility.

The mediators used in this analysis (tokens, compliance, and duration) are all post-treatment variables that themselves respond to the prompt framing. Because these variables change in response to the treatment, adjusting for them cannot isolate a clean causal direct effect, and we do not interpret the adjusted coefficients as proving a mechanism. The analysis is nonetheless useful for system design when stated in associational terms. If a pass-rate contrast is attenuated after adjustment for measured behavioral variables, those behaviors are associated with part of the observed contrast, a pattern consistent with a parser-compatibility or verbosity pathway. Conversely, if a contrast remains after adjustment, the measured behaviors account for only part of it, and the remainder reflects unmeasured aspects of the model--evaluator interaction that this study cannot identify. In both cases, the interpretation relies on observable outputs rather than hidden reasoning processes.

\section{Dominant-Model Sensitivity}

The dominant-model analysis comprises a two-level robustness check. First, it examines how the pooled pass rates vary when the most powerful model within the group is excluded from the analysis. In the second and more critical step, it investigates whether the prompt effects themselves are distorted. This re-estimation takes place for each main tone, role, and format contrast, with and without the dominant model, relying on both the LPM percentage-point scale and the mixed-effects logistic scale. In this way, it becomes possible to assess the direction and magnitude of prompt effects across the panels.

In small model panels, this leave-one-dominant-model approach is particularly significant. When one model has an extraordinarily high pass rate, pooled prompt effects may simply reflect how that specific model responds to the prompts. Therefore, reporting only the pooled effect of all models would create a misleadingly generalized result. The sensitivity analysis seeks to determine whether the same practical outcome still exists once this dominant model is extracted from the data. If it no longer exists, the thesis will directly state this dependence rather than smooth it out.

A structural limitation applies to the API group. Because there are only two cloud models, excluding DeepSeek V4 Flash leaves a single-model panel (comprising only Gemini Flash), meaning the leave-one-out estimates describe a single model rather than a model average. We report these estimates because they directly expose how much of each pooled API effect is attributable to the dominant model.

All of these individual analyses serve as the basis for the conservative evidentiary standard employed in this thesis. The strongest evidence is provided by patterns across tasks and model groups, similarity across summaries, and the impact on system behavior. Conversely, the weakest evidence is provided by a single p-value, pooled averages dominated by one model, or system behavior that does not align with correctness. This evidentiary standard reflects the goal of this thesis: to view prompt framing as an engineering control variable, rather than a universal benchmark.
